\documentclass[11pt]{article}

\usepackage[margin=1in]{geometry}
\usepackage{amsmath,amssymb,amsthm}
\usepackage{booktabs}
\usepackage{enumitem}
\usepackage[colorlinks=true,citecolor=blue,linkcolor=blue,urlcolor=blue]{hyperref}
\usepackage[round,authoryear]{natbib}
\usepackage{times}
\usepackage{xcolor}

\newcommand{\eps}{\varepsilon}
\newcommand{\Real}{\mathbb{R}}

% Simple, dependency-free callout blocks: an indented quote with a bold
% running head and a vertical rule down the left margin, built only from
% core LaTeX primitives so the document does not depend on a fragile
% framing-package chain.
\newenvironment{intuition}
  {\par\addvspace{8pt}\begin{quote}\noindent\textbf{\textcolor{blue!40!black}{Intuition.}}\ }
  {\end{quote}\par\addvspace{8pt}}

\newenvironment{worked}
  {\par\addvspace{8pt}\begin{quote}\noindent\textbf{\textcolor{black!50!white}{Worked example.}}\ }
  {\end{quote}\par\addvspace{8pt}}

\title{A Plain-Language, Technical Guide to FlowSurv-AFT:\\
What This Research Does, What Came Before It, and Why It Is Different}
\author{}
\date{}

\begin{document}
\maketitle

\begin{abstract}
\noindent
This document is written for a reader who already understands statistics,
machine learning, and deep learning in general, but who has never worked
with \emph{survival analysis} specifically. It builds up the subject from
first principles -- what survival data looks like, why it cannot be
handled by ordinary regression, and what mathematical objects (the
survival function, the hazard function, the censored likelihood) the
field uses instead -- and then walks, in order, through every major family
of survival model that has been proposed, explaining in plain language
\emph{how each one works mechanically}, not just what it is called. With
that foundation in place, it explains this research's method,
\emph{FlowSurv-AFT}, in full mathematical detail, and finally explains
precisely why it differs from every prior approach, including the two
closest competing methods. There is no page limit on this document and no
attempt to compress the explanation; the goal is that a careful reader
finishes it able to explain the whole research programme to someone else.
\end{abstract}

\tableofcontents

\newpage

\section{Why survival analysis needs its own toolkit}
\label{sec:why}

Suppose you want to predict how long a patient will live after a cancer
diagnosis, how long a machine part will run before it breaks, or how long
a customer will stay subscribed to a service before cancelling. At first
glance this looks like an ordinary regression problem: you have a
continuous outcome (a time) and some covariates (patient age, tumour
size; machine load, material batch; customer usage pattern), so why not
just fit a neural network or a linear model to predict the time directly?

The reason is a single, unavoidable fact about how this kind of data is
collected: \textbf{the study ends before every subject's event has
happened.} If you follow 1{,}000 patients for five years, some of them
will die during the study -- for those, you know the exact time. But many
will still be alive when the study ends, and some will drop out partway
through for unrelated reasons (they moved away, they withdrew consent).
For all of those subjects, you do \emph{not} know their true survival
time. All you know is a \emph{lower bound}: "this patient was alive for
at least 3.2 years." This phenomenon is called \textbf{censoring}, and it
is the single fact that makes survival analysis a distinct field rather
than a special case of regression.

\begin{intuition}
If you simply threw away everyone who was still alive at the end of the
study and fit a regression model only to the patients who died, you would
get a badly biased answer: you would be systematically excluding your
\emph{longest}-surviving patients, so your model would think everyone
dies sooner than they really do. If instead you treated "study ended at
5 years, patient still alive" as though it meant "patient's survival time
is exactly 5 years," you would also be wrong, in the opposite direction --
you would be throwing away the information that this patient survived
\emph{at least} 5 years and might have survived much longer. Survival
analysis exists to use exactly the right amount of information contained
in a censored observation -- no more, no less.
\end{intuition}

This single complication -- that some outcomes are only partially
observed -- forces every survival model to be built around a specific
mathematical machinery (Section~\ref{sec:math}) that ordinary regression
and classification models do not need. Once you understand that
machinery, every model discussed in this document, including this
research's own model, is a variation on the same underlying idea: specify
a probability distribution for the event time, and fit it in a way that
correctly uses both the fully-observed event times \emph{and} the
partially-observed censored ones.

\subsection{Types of censoring}
\label{sec:censoring-types}

There is more than one way an observation can be incomplete, and the type
matters for how the data should be modelled.

\begin{itemize}
  \item \textbf{Right censoring} is the case described above: we know the
    subject survived at least until some time $C$, but not what happened
    after. This is overwhelmingly the most common case in practice (a
    study ends; a patient is lost to follow-up) and is the primary focus
    of this research. Right censoring itself comes in a few flavours that
    matter for how it is generated in a simulation study:
    \begin{itemize}
      \item \emph{Type~I censoring}: every subject has the \emph{same}
        fixed censoring time (e.g., "the study runs for exactly 5
        years, and everyone who hasn't had the event by then is
        censored at 5 years").
      \item \emph{Type~III (random) censoring}: each subject has their
        \emph{own}, independently drawn censoring time (e.g., people
        enrol in the study at different calendar dates, or drop out for
        unrelated personal reasons at random times).
    \end{itemize}
  \item \textbf{Left censoring} is the mirror case: we know the event
    happened \emph{before} some observed time, but not exactly when
    (e.g., a screening test detects a disease that could have started
    at any point before the test date).
  \item \textbf{Interval censoring} is when we only know the event
    happened somewhere between two observation times (e.g., a patient is
    disease-free at one check-up and has the disease at the next, but the
    exact onset time in between is unknown).
\end{itemize}

This research's primary setting is right censoring, under both Type~I and
Type~III mechanisms, with interval censoring treated as a secondary
experiment. A further, harder complication -- \emph{dependent} or
\emph{informative} censoring, where the reason a subject is censored is
itself related to their unobserved event time (e.g., patients who are
about to die are more likely to drop out of a study) -- is explicitly
\emph{not} addressed by this research; it is a separate, harder problem
addressed by other work \citep{gharari2023copula}, and every model
discussed in this document, including FlowSurv-AFT, assumes censoring is
independent of the event time given the covariates.

\section{The mathematical language of survival analysis}
\label{sec:math}

This section introduces the handful of mathematical objects that every
survival model, from the oldest (1950s actuarial tables) to the newest
(this research's own deep flow model), is built out of. If you already
know what a survival function and a hazard function are and how they
relate, you can skim this section; if not, read it carefully, because
every later section assumes it.

\subsection{The random variable, and three equivalent ways to describe its distribution}

Let $T$ be the (true, possibly unobserved) time until the event of
interest, a non-negative continuous random variable. As with any
continuous random variable, $T$ has a cumulative distribution function
$F(t) = \Pr(T \le t)$ and a probability density function $f(t) =
F'(t)$. Survival analysis, however, almost always works with the
\emph{complement} of $F$, called the \textbf{survival function}:
\begin{equation}
  S(t) \;=\; \Pr(T > t) \;=\; 1 - F(t).
  \label{eq:survival-def}
\end{equation}
$S(t)$ answers the question "what fraction of the population is still
event-free after time $t$?" -- it starts at $S(0) = 1$ (everyone is
event-free at time zero) and decreases monotonically toward $0$ as $t \to
\infty$ (eventually, everyone experiences the event, assuming $T$ is not
defective). A survival curve that stays close to 1 for a long time and
then drops describes a population where the event is rare until some
point and then common; a survival curve that decays smoothly and slowly
describes constant risk over time.

The third, and for this research the most important, way to describe the
distribution of $T$ is the \textbf{hazard function}:
\begin{equation}
  h(t) \;=\; \lim_{\Delta t \to 0} \frac{\Pr(t \le T < t + \Delta t \mid T \ge t)}{\Delta t}.
  \label{eq:hazard-def}
\end{equation}
This is an \emph{instantaneous rate}: given that a subject has survived
up to time $t$, $h(t)$ is (informally) the rate at which the event is
about to happen right now, per unit time. It is not a probability -- it
can exceed 1 -- and it is best read as "risk per unit time, right now,
conditional on having made it this far."

\begin{intuition}
The hazard function is the single most informative object in survival
analysis because it is the one most directly tied to the physical or
biological \emph{mechanism} generating the data. A rising hazard means
the risk of the event is accelerating over time (wear-out, ageing); a
falling hazard means early survivors are increasingly "safe" (infant
mortality, burn-in failures); a hazard that falls and then rises again
-- a \textbf{bathtub} shape -- describes a population with high early
risk (surgical complications, manufacturing defects), a stable middle
period, and rising late risk (ageing, material fatigue); a hazard with
more than one bump (\textbf{multimodal}) describes a population that is
secretly a mixture of sub-groups with different risk timelines pooled
under one label. Every model discussed in this document should be judged,
first, by which of these shapes it is even \emph{capable} of representing
-- a question about the model's functional form, separate from how much
data is available to fit it.
\end{intuition}

\subsection{How the three objects relate to each other}

$f$, $S$, and $h$ are not independent choices -- any one of them
determines the other two. First, by definition, $f(t) = -S'(t)$ (the
density is minus the rate of decay of the survival curve). Second, a short
derivation from Eq.~\eqref{eq:hazard-def} gives the hazard directly in
terms of $f$ and $S$:
\begin{equation}
  h(t) \;=\; \frac{f(t)}{S(t)}.
  \label{eq:hazard-ratio}
\end{equation}
This says exactly what the intuition box above suggested: the hazard is
the density of "the event happens right now" \emph{rescaled} by the
probability of still being at risk at all. Finally, integrating
Eq.~\eqref{eq:hazard-ratio} (using $f = -S'$) gives the single most
useful identity in the whole field, connecting the hazard to the survival
function directly:
\begin{equation}
  S(t) \;=\; \exp\!\left(-\int_0^t h(u)\,du\right) \;=\; \exp\bigl(-H(t)\bigr),
  \qquad H(t) := \int_0^t h(u)\,du,
  \label{eq:survival-from-hazard}
\end{equation}
where $H(t)$ is called the \textbf{cumulative hazard}. Every survival
model, no matter how it is parameterized, is implicitly (or explicitly)
choosing a functional form for one of $\{f, S, h, H\}$ and deriving the
other three from Eqs.~\eqref{eq:hazard-ratio}--\eqref{eq:survival-from-hazard}.
Whether those derivations can be done \emph{in closed form} -- with a
formula you can write down and evaluate directly -- or only
\emph{numerically} -- by running an algorithm (an ODE solver, a root
finder, a Monte Carlo average) every single time you want a number out of
the model -- turns out to be one of the central technical distinctions
this document uses to compare methods.

\subsection{A worked example: the Weibull distribution}

\begin{worked}
The Weibull distribution is the classical workhorse of survival analysis
and reliability engineering because it has a simple, interpretable hazard
function. With shape parameter $\rho > 0$ and scale parameter $\lambda >
0$,
\[
  h(t) = \lambda \rho (\lambda t)^{\rho - 1}, \qquad
  S(t) = \exp\bigl(-(\lambda t)^\rho\bigr), \qquad
  f(t) = h(t)\,S(t).
\]
Look at the hazard formula: if $\rho > 1$, $h(t)$ is increasing in $t$
(wear-out); if $\rho < 1$, $h(t)$ is decreasing (infant mortality); if
$\rho = 1$, $h(t) = \lambda$ is constant (this special case is the
\emph{exponential} distribution, "memoryless" constant risk). Crucially,
for \emph{no} value of $\rho$ does $h(t)$ turn over and change direction:
a single Weibull hazard is monotone for every $\rho \ne 1$, full stop.
This is not a limitation of how much data you have -- it is baked into
the functional form. A bathtub-shaped hazard, or a hazard with two
bumps, is a shape a single Weibull distribution can \emph{never}
represent, no matter the sample size. This single fact -- that a fixed
parametric family imposes a hard ceiling on which hazard shapes are even
representable -- is the motivating problem of this entire research
programme, and it recurs, in one guise or another, in Sections
\ref{sec:classical}--\ref{sec:generative} below.
\end{worked}

\subsection{The censored likelihood: using exactly the right information}
\label{sec:censored-likelihood}

Now return to the central complication from Section~\ref{sec:why}: for
each subject $i$ we do not observe $T_i$ directly. We observe a pair
$(\tilde t_i, \delta_i)$, where
\[
  \tilde t_i = \min(T_i, C_i), \qquad
  \delta_i = \mathbb{1}\{T_i \le C_i\} \;\; \text{(1 if the event was observed, 0 if censored)},
\]
with $C_i$ the (independent) censoring time. How should a model be fit to
data like this? The key idea -- and the single most important
construction in the entire field -- is to write down the \emph{correct}
likelihood contribution for each subject given what was actually
observed, rather than pretending the data is uncensored.

\begin{itemize}
  \item If $\delta_i = 1$ (event observed at $\tilde t_i$): we know
    \emph{exactly} when the event happened. The likelihood contribution
    is the density evaluated at that time, $f(\tilde t_i \mid x_i)$ --
    exactly as in ordinary maximum-likelihood regression.
  \item If $\delta_i = 0$ (censored at $\tilde t_i$): we know only that
    the event has \emph{not yet happened} by time $\tilde t_i$. The
    likelihood contribution is not a density value at all, but a
    \emph{probability}: $\Pr(T_i > \tilde t_i \mid x_i) = S(\tilde t_i
    \mid x_i)$.
\end{itemize}

Combining both cases into one formula (using $\delta_i$ as a $0/1$
switch), the contribution of subject $i$ to the likelihood is
\begin{equation}
  L_i \;=\; f(\tilde t_i \mid x_i)^{\delta_i}\; S(\tilde t_i \mid x_i)^{1 - \delta_i},
  \label{eq:censored-likelihood}
\end{equation}
and the full log-likelihood to be maximized over the whole dataset is
\begin{equation}
  \ell \;=\; \sum_{i=1}^n \Bigl[\delta_i \log f(\tilde t_i \mid x_i) + (1-\delta_i)\log S(\tilde t_i \mid x_i)\Bigr].
  \label{eq:censored-loglik}
\end{equation}
This is the \textbf{exact right-censored likelihood}, and it is the
estimation criterion this research's model is fit with. It is worth
dwelling on why it is called \emph{exact}: Eq.~\eqref{eq:censored-loglik}
is a direct, unmodified application of the maximum-likelihood principle
to the data that was actually observed -- it is not an approximation, a
re-weighting scheme, or a workaround. As Section~\ref{sec:deepaft}
explains in detail, a great deal of the deep-learning survival literature
does \emph{not} fit models this way, instead substituting an
approximate or indirect estimation criterion (rank losses, inverse
weighting, iterative imputation) -- usually because the model's chosen
functional form for $f$ or $S$ makes Eq.~\eqref{eq:censored-loglik}
inconvenient or impossible to evaluate directly. Whether a model can be
trained by the honest formula in Eq.~\eqref{eq:censored-loglik}, or
needs a substitute, is the second central technical distinction this
document uses to compare methods (the first being closed-form versus
numerical evaluation, from the previous subsection).

\section{How survival models have traditionally been evaluated}
\label{sec:evaluation}

Before touring the models themselves, it helps to know how their quality
is measured, because several of the historically popular metrics turn
out to be more subtle -- and less trustworthy as a basis for comparing
methods -- than they first appear. This section is necessary background
for Section~\ref{sec:generative} onward, where evaluation choices become
part of the story.

\subsection{Discrimination: can the model rank subjects correctly?}

The most widely reported metric is a \textbf{concordance index}
(C-index), which asks: if you take two subjects at random (such that we
can tell who had the earlier true event time), does the model's
predicted risk score correctly rank them? Concretely, for every
comparable pair $(i, j)$ with $\tilde t_i < \tilde t_j$ and $\delta_i =
1$, check whether subject $i$'s predicted risk is higher than subject
$j$'s; the C-index is the fraction of comparable pairs the model gets
right (0.5 = no better than a coin flip, 1.0 = perfect ranking).

\begin{worked}
Suppose the model predicts risk scores $r_A = 0.8$, $r_B = 0.3$ for two
patients, and patient $A$ truly died before patient $B$ (with $A$'s death
observed, not censored, so the pair is comparable). The model correctly
ranked $A$ as riskier, so this pair counts as \emph{concordant}. If
instead $r_A = 0.2 < r_B = 0.3$, the pair would be \emph{discordant}. If
patient $B$ were censored \emph{before} patient $A$'s event, we could not
tell who truly had the higher risk, so the pair is simply excluded from
the calculation. The C-index is (concordant pairs) $/$ (concordant $+$
discordant pairs).
\end{worked}

Two subtleties matter a great deal for this research. First, the
"classic" concordance index (Harrell's C) is known to be biased under
heavy censoring, because pairs where the earlier subject is censored are
silently dropped, and dropping is not random with respect to risk; Uno's
C \citep{uno2011} corrects this with an inverse-probability-of-censoring
weighting scheme. Second, and more importantly, \citet{sonabend2022chacking}
points out a practice the field calls \textbf{"C-hacking"}: because there
are several competing definitions of concordance (Harrell's, Uno's,
time-dependent variants, and others), a researcher can always find
\emph{some} variant under which their method looks best and report only
that one. The discipline against this is to pre-specify, before running
any experiment, exactly one concordance variant that will be reported --
which is what this research does (Uno's C, chosen in advance).

\subsection{Overall accuracy: the Brier score}

The concordance index only measures \emph{ranking}, not whether the
model's predicted \emph{probabilities} are numerically close to the
truth. The Integrated Brier Score \citep{graf1999ibs} fixes this by
computing, at a set of time horizons $t$, the squared error between the
model's predicted survival probability $\hat S(t \mid x_i)$ and the
actual observed outcome $\mathbb{1}\{\tilde t_i > t\}$, averaged over
subjects and integrated over the time horizons of interest -- with an
inverse-probability-of-censoring correction \citep{gerds2006} so that
censored subjects still contribute a valid, unbiased term to the
average rather than being silently dropped as in a naive squared-error
calculation.

\subsection{Calibration: are the predicted \emph{distributions} correct?}

This is the metric family that matters most for the claims this research
makes, and it is the one an ML/DL reader is least likely to have
encountered before, because it has no exact analogue in ordinary
classification or regression evaluation. A model can rank subjects
perfectly (high C-index) while still being badly \textbf{miscalibrated}
-- systematically predicting survival probabilities that are too high or
too low in absolute terms, even though the relative ordering is right.

\textbf{D-calibration} \citep{haider2020} tests this directly: take every
subject's predicted survival probability \emph{at their own observed
time}, $\hat S(\tilde t_i \mid x_i)$ (for an event-observed subject, this
should behave like a draw from $\mathrm{Uniform}(0,1)$ if the model is
correct -- a standard probability-integral-transform argument), bin these
values into deciles, and run a chi-squared test for whether the bins are
uniformly populated as they should be under a correctly specified model.
The \textbf{integrated calibration index} (ICI, \citealp{austin2020ici})
instead plots predicted probability against observed (smoothed) outcome
frequency at a fixed time horizon and reports the average absolute gap
between the two -- the survival-analysis analogue of a reliability
diagram, familiar from classifier calibration, but adapted to handle
censored outcomes correctly.

\subsection{The evidence that reframes the whole comparison}

The single most important piece of evaluation-methodology evidence this
research relies on is empirical rather than definitional. In the largest
neutral comparison conducted to date -- 19 different survival methods,
spanning classical, machine-learning, and deep approaches, evaluated
across 34 real datasets under one common, fair protocol --
\citet{burk2024neutral} find that \emph{no method, including every deep
learning method surveyed, significantly outperforms plain Cox regression
in discrimination on standard, low-dimensional tabular survival data}.

\begin{intuition}
Why does this matter for how to read the rest of this document? It means
that if this research (or any deep survival model) tried to justify
itself by reporting a better C-index than Cox regression, the honest
prior expectation, based on the best available neutral evidence, is that
no such improvement is really there to be found -- any apparent gain
would likely be noise, or an artefact of an unfair comparison protocol.
The scientifically defensible place for a flexible, deep survival model
to add value is therefore \emph{not} discrimination, but
\textbf{calibration and distributional shape fidelity} -- whether the
model's predicted probability distribution over event times is
numerically correct and captures the true hazard shape, not merely
whether it ranks subjects in the right order. This research's central
claims, throughout, are built around calibration and hazard-shape
recovery for exactly this reason, and it deliberately avoids leading with
a concordance claim.
\end{intuition}

\section{Classical survival models}
\label{sec:classical}

\subsection{The Kaplan--Meier curve: the simplest possible starting point}

Before any covariates enter the picture, the most basic survival
estimator is the \textbf{Kaplan--Meier} curve: a step-function estimate
of $S(t)$ built directly from the censored data, with no parametric
assumptions and no covariates at all. It works by multiplying together,
at each observed event time, the fraction of at-risk subjects who
survived past that instant, correctly "skipping over" censored subjects
by keeping them in the risk set up until the moment they are censored,
then removing them without counting them as an event. It has no closed
form and does not condition on covariates, so it is not itself a
predictive model, but it is the estimator every more sophisticated method
in this document is implicitly trying to improve upon.

\subsection{The Cox proportional hazards model}
\label{sec:cox}

The Cox model \citep{cox1972} is the field's default tool for
covariate-adjusted analysis. Its defining assumption is
\textbf{proportional hazards}:
\begin{equation}
  h(t \mid x) \;=\; h_0(t)\,\exp(\beta^\top x),
  \label{eq:cox}
\end{equation}
where $h_0(t)$ is an unspecified \textbf{baseline hazard} (the hazard for
a subject with $x = 0$) and $\beta$ are covariate coefficients. The word
"proportional" describes the key structural consequence: the
\emph{ratio} of hazards between any two covariate profiles,
$h(t\mid x_a)/h(t\mid x_b) = \exp\{\beta^\top(x_a - x_b)\}$, does not
depend on $t$ at all -- one profile is always some fixed multiple riskier
than the other, at every point in time. This is a strong, testable
assumption, and it fails whenever a covariate's effect genuinely changes
over time (a classic example: a surgical intervention that reduces
short-term risk but has no long-term effect, or even reverses -- a
"crossing hazards" scenario).

The reason the Cox model is called \textbf{semiparametric} is that
$\beta$ can be estimated by maximizing a \textbf{partial likelihood} that
depends only on the \emph{ranking} of event times among at-risk subjects,
without ever specifying $h_0(t)$ at all. This is elegant and one reason
for the model's enduring popularity, but it has a direct cost relevant to
this document: because $h_0(t)$ is never estimated as a smooth function
(only recoverable afterward as a non-smooth Breslow-type step estimate),
the Cox model produces \emph{no smooth density $f(t\mid x)$, no smooth
hazard curve, and no closed-form quantile function} -- exactly the
objects this research's model is built to deliver in closed form.

\subsection{The accelerated failure time (AFT) model}
\label{sec:aft}

The AFT model \citep{wei1992aft} takes an entirely different approach: it
specifies a complete parametric distribution for $\log T$ directly,
\begin{equation}
  \log T \;=\; \beta^\top x + \sigma\,\eps, \qquad \eps \sim F_0,
  \label{eq:aft}
\end{equation}
where $F_0$ is a fixed, standard reference distribution (common choices:
the extreme-value distribution, which makes $T$ Weibull-distributed; the
normal distribution, which makes $T$ log-normal; the logistic
distribution, which makes $T$ log-logistic), $\beta^\top x$ is a linear
"location" term, and $\sigma$ is a scale parameter. Because $F_0$ is
fully specified, \emph{everything} -- the density, the survival function,
the hazard, and every quantile of $T$ -- is available in closed form as
soon as $\beta$ and $\sigma$ are estimated (by directly maximizing
Eq.~\eqref{eq:censored-loglik}, no partial-likelihood trick needed).

\begin{intuition}
The name "accelerated failure time" comes directly from what happens to
the whole distribution when you change $x$. Compare two covariate
profiles $x_a$ and $x_b$. From Eq.~\eqref{eq:aft}, $\log T_a - \log T_b =
\beta^\top(x_a - x_b)$ when the same underlying $\eps$ is used for both
(a coupling argument), so
\begin{equation}
  T_a \;=\; T_b \cdot \exp\{\beta^\top(x_a - x_b)\} \;=\; T_b \cdot \mathrm{TR},
  \qquad \mathrm{TR} := \exp\{\beta^\top(x_a - x_b)\}.
  \label{eq:time-ratio}
\end{equation}
Every event time for profile $a$ is exactly $\mathrm{TR}$ times the
corresponding event time for profile $b$ -- \emph{not just the median or
the mean, but every quantile of the distribution simultaneously.} If
$\mathrm{TR} = 2$, profile $a$'s entire lifetime distribution is stretched
to twice the length of profile $b$'s; if $\mathrm{TR} = 0.5$, it is
compressed to half. This \textbf{time-ratio} or \textbf{acceleration
factor} is, for many practitioners -- especially in reliability
engineering, where "expected lifetime is doubled" feeds directly into
warranty and maintenance planning -- a substantially more intuitive
quantity to communicate than a Cox hazard ratio, which only describes an
instantaneous relative rate, not how the whole distribution moves.
\end{intuition}

The AFT's cost is exactly the flip side of the Cox model's: by fixing
$F_0$ to a specific standard distribution, the model inherits that
distribution's hazard-shape restrictions. As the worked Weibull example
in Section~\ref{sec:math} showed, a Weibull AFT can never represent a
bathtub hazard; a log-normal AFT is likewise restricted to a single
hump-shaped hazard. This is the central tension this entire research
programme resolves: \emph{can we keep the AFT's closed-form tractability
and time-ratio interpretability, while removing the restriction that
$F_0$ must be a simple, fixed, named distribution?}

\subsection{Flexible parametric (Royston--Parmar) models}
A middle ground between Cox's total shape-freedom (but no smooth density)
and the AFT's shape restriction (but full tractability) is
\citet{royston2002}'s flexible parametric model, which represents the log
cumulative hazard as a \textbf{restricted cubic spline} -- a flexible,
piecewise-polynomial curve, smoothly joined at a handful of "knot"
locations -- in $\log t$, rather than committing to a named parametric
family:
\[
  \log H(t \mid x) \;=\; s(\log t;\, \gamma) + \beta^\top x.
\]
This is the strongest classical shape-flexible baseline available (and is
one of the baseline models this research compares against), but because
the spline is a free-form nonparametric object rather than a
probability-generating transformation applied to a single latent
variable, it does not, without further restriction, carry the
location-scale AFT interpretation of Eq.~\eqref{eq:aft}: there is no
single time-ratio quantity that summarizes the covariate effect the way
$\mathrm{TR}$ in Eq.~\eqref{eq:time-ratio} does.

\subsection{Random Survival Forests}
Random Survival Forests \citep{ishwaran2008rsf} extend the idea of a
decision-tree ensemble (familiar from ordinary classification/regression
forests) to censored survival data, splitting on covariates to group
similar subjects and averaging a Nelson--Aalen-type cumulative hazard
estimate within each leaf. Like Cox regression, this handles nonlinearity
and interactions without committing to a hazard shape, but -- again like
Cox regression -- it produces no smooth conditional density, only a
step-function-like ensemble curve, and its calibration properties are
comparatively weak.

\section{Deep learning survival models that keep a classical structure}
\label{sec:dl}

The next three families all apply neural networks to survival data, but
each one does so by extending -- and thereby inheriting the limitations
of -- one of the classical structures from Section~\ref{sec:classical}.

\subsection{Deep Cox models}
DeepSurv \citep{katzman2018deepsurv} is the most direct: it replaces the
linear predictor $\beta^\top x$ inside the Cox partial likelihood
(Eq.~\eqref{eq:cox}) with a neural network's output, $g_\theta(x)$,
gaining the ability to model nonlinear covariate effects and
interactions, but keeping everything else about the Cox model unchanged
-- proportional hazards, an unspecified baseline hazard, partial (not
full) likelihood, no smooth density. Cox-Time \citep{kvamme2019coxtime}
extends this further by letting the network also depend on time,
$g_\theta(x, t)$, which lets the model escape the strict proportionality
assumption (hazard \emph{ratios} can now vary in time), but the
estimator still factors through a shared, non-parametrically estimated
baseline hazard, so the output remains a step-function survival curve
rather than a smooth density.

\subsection{Discrete-time models: DeepHit and friends}
\label{sec:deephit}
A different family sidesteps the Cox baseline-hazard machinery entirely
by discretizing time onto a grid of bins $[0, \tau_1), [\tau_1, \tau_2),
\ldots$, and having the network directly output a probability for
"the event happens in this bin," for every bin. DeepHit
\citep{lee2018deephit} is the best-known example: the network outputs a
softmax distribution over the bins (a genuine, valid probability mass
function that sums to 1, exactly the mechanism a classification network
would use over class labels), trained with a loss that combines a
likelihood-like term with a ranking term. This is a clever way to get a
proper, trainable probability distribution over event time without
needing to assume any parametric shape, and it can natively handle
\emph{competing risks} (more than one possible type of event) by having
one softmax per risk type.

\begin{intuition}
The cost of this trick is what we call the \textbf{discretization horn}
throughout this document: because the model's output is a probability
\emph{per bin}, not a density function of a continuous variable, the
implied survival curve is a step function, jumping only at bin
boundaries; there is no density or hazard value defined at an arbitrary
time $t$ that falls inside a bin (only a bin-level probability); and any
quantile the model reports (including the median survival time used to
rank subjects for concordance) is only resolved to the width of a bin.
Choosing the grid is therefore not a cosmetic detail but a real
bias--variance trade-off baked directly into the estimator: a coarse
grid with few, wide bins discards real temporal resolution the data may
actually support; a fine grid with many, narrow bins multiplies the
number of parameters to be estimated and can make training unstable
under heavy censoring, where many bins may receive very few observed
events at all.
\end{intuition}

Related discrete-time and interval-hazard models follow the same basic
pattern with variations: DRSA \citep{ren2019drsa} predicts the same kind
of per-bin probabilities autoregressively (one bin at a time, each
conditioned on the previous); PC-Hazard \citep{kvamme2021pchazard}
instead models a piecewise-\emph{constant} hazard within each bin, which
admits a cleaner (piecewise) likelihood derivation; SurvTRACE
\citep{wang2022survtrace} keeps the same discrete-hazard output but
replaces the covariate encoder with a transformer architecture. All four
inherit the same fundamental limitation described in the intuition box
above.

\subsection{Continuous-time flexible models}
A smaller family tries to keep time genuinely continuous without
discretizing it. SODEN \citep{tang2022soden} generalizes the DeepSurv /
DeepHit idea through a system of ordinary differential equations (a
\textbf{neural ODE}) whose numerical solution implicitly defines the
hazard function -- gaining real flexibility, but at the cost of a density
that is never available in closed form and must be recomputed by solving
the ODE system every single time the model is evaluated. SuMo-net
\citep{rindt2022sumonet} instead has a neural network directly
parameterize $S(t \mid x)$, constrained by its architecture to always be
monotonically decreasing in $t$, trained with a proper scoring rule
rather than the exact likelihood; this gives a genuinely smooth survival
curve, but the corresponding density and hazard exist only by
numerically \emph{differentiating} the learned survival curve, which is
not exact and tends to amplify estimation noise exactly in the tails
where questions about bathtub or multimodal hazard shapes are decided.

\section{Deep accelerated failure time models: a flexible location, but no distribution}
\label{sec:deepaft}

The natural next idea, following Section~\ref{sec:aft}, is to keep the
AFT's Eq.~\eqref{eq:aft} but replace the linear location term
$\beta^\top x$ with a neural network, $\mu_\theta(x)$. Several papers do
exactly this -- and this section explains, in mechanical detail, the
specific estimation trade-off every one of them makes.

\subsection{Why "just add a neural network" breaks the exact likelihood}
The AFT's closed-form tractability in Section~\ref{sec:aft} came from a
specific fact: $F_0$ was a \emph{named, fixed} distribution with a known
density formula. If you keep $F_0$ fixed but let $\mu_\theta(x)$ be an
arbitrary neural network, the likelihood in
Eq.~\eqref{eq:censored-loglik} is still perfectly well-defined and
computable -- this is, in fact, exactly what this research's own
$\mu(x)$ head does (Section~\ref{sec:model}). The papers discussed in
this section, however, additionally leave the \emph{shape} of $F_0$
unspecified (not just the location), which is precisely what removes the
closed-form density needed to evaluate Eq.~\eqref{eq:censored-loglik}
directly -- and it is at that point that each of these methods reaches
for a substitute estimation criterion instead of the exact likelihood.

\subsection{Buckley--James estimation and inverse-probability weighting}
deepAFT \citep{norman2024deepaft} places a deep network on $\mu(x)$
without committing to any $F_0$ at all, and estimates it using one of
several classical semiparametric AFT-fitting techniques originally
developed for the linear model: \textbf{Buckley--James} iterative
imputation, which repeatedly replaces each censored observation's
"missing" residual with its conditional expectation under the current
model fit and re-fits, iterating to convergence; or
\textbf{inverse-probability-of-censoring weighting (IPCW)}, which
reweights each observation by the inverse of its estimated probability of
\emph{not} having been censored by its observed time, so that
uncensored observations "stand in" for the censored ones they are
weighted to represent.

\begin{intuition}
Both techniques are well-established and valid in the classical
(linear-$\mu$) AFT setting, but both share a structural weak point that
matters directly for this research's design: their reliability degrades
as the censoring proportion grows. Buckley--James imputation becomes
increasingly reliant on extrapolating from a shrinking pool of
fully-observed residuals as more of the data is censored. IPCW weights
are the \emph{reciprocal} of an estimated survival probability for
censoring; as that probability approaches zero (which becomes more
common precisely when overall censoring is heavy), the weights blow up,
and a small number of observations can dominate the entire fitted
estimate. Because this research's simulation study explicitly sweeps
censoring proportion from 0\% up to 80\%, this is exactly the regime
where such estimators are expected, on general statistical grounds, to
become least reliable -- which is part of the motivation for insisting
on the exact likelihood (Eq.~\eqref{eq:censored-loglik}) instead.
Whatever technique is used, the resulting deepAFT survival curve is a
Kaplan--Meier-rescaled step function, not a smooth conditional density.
\end{intuition}

Two further variants follow the same likelihood-free pattern with a
different neural architecture or training signal. DeepR-AFT
\citep{kim2022deepraft} trains $\mu_\theta(x)$ using a
\textbf{Gehan-type rank loss} -- a pairwise loss that only rewards
getting the \emph{relative order} of predicted times right (structurally
similar in spirit to the concordance idea from
Section~\ref{sec:evaluation}), which supplies a valid training signal for
ranking but carries no information at all about the scale or shape of
the predicted distribution, so no calibrated density can be extracted
from the fitted model. KAN-AFT \citep{francis2025kanaft} substitutes a
Kolmogorov--Arnold network (a recent neural architecture that
represents functions as sums of learned univariate splines rather than
matrix-multiplication layers, offering a claimed advantage in symbolic
interpretability of the fitted $\mu(x)$) for $\mu_\theta(x)$, but keeps
the same likelihood-free estimation as deepAFT. Two further recent
entries, GRAFT \citep{graft2026} and DART \citep{dart2023}, follow the
identical pattern -- a flexible AFT-style location network trained with a
concordance-aligned rank loss (with, in GRAFT's case, a post-hoc
calibration correction bolted on afterward) -- and so occupy exactly the
same position in this taxonomy.

\subsection{Keeping the likelihood by fixing the shape instead}
Two papers make the opposite trade: they keep the exact censored
likelihood, but pay for it by fixing $F_0$ to a single, simple shape.
Neural Conditional Event Time models \citep{engelhard2020neuralcet} --
"Neural-CET" -- place neural networks on \emph{both} the location $\mu(x)$
\emph{and} the scale $\sigma(x)$ of Eq.~\eqref{eq:aft}, and fit the model
by the exact likelihood of Eq.~\eqref{eq:censored-loglik} -- but with
$F_0$ fixed to the standard normal distribution. Structurally, this
means Neural-CET \emph{is} a log-normal AFT model with a neural covariate
map: the location and scale can be arbitrarily flexible functions of
$x$, but the \emph{shape} of the residual distribution around that
location and scale is permanently fixed to the single, symmetric,
unimodal hump of the normal distribution -- exactly as shape-restricted,
in principle, as the classical log-normal AFT of Section~\ref{sec:aft}.
This is, in fact, the single closest prior model to the present
research's own architecture, and this research includes an ablation
model, called \emph{FlowSurv-Gauss}, that is constructed to be exactly
this: the same neural $\mu(x)/\sigma(x)$ encoder as the full model, but
with the flexible residual transform switched off and the base
distribution fixed to Gaussian -- so that comparing the full model
against this ablation isolates, as precisely as an experiment can, how
much of any performance improvement comes specifically from freeing the
residual shape rather than from any other architectural difference. DKAFT
\citep{wu2021dkaft} makes a similar move using a Gaussian-process output
layer in place of a plain neural network, but is likewise log-normal in
effect. \citet{ranganath2016deepsurv}'s deep exponential family
parameterizes Weibull-distributed event times through a deep generative
model -- again, the exact likelihood is retained, but the fixed Weibull
shape's monotone-hazard restriction (Section~\ref{sec:math}'s worked
example) is inherited unchanged, however deep the covariate network
becomes.

\subsection{The pattern, stated once, clearly}
Putting the two halves of this section side by side reveals a single,
consistent pattern, and it is the pattern this research is built to
break: \emph{every deep AFT model reviewed above either (a) keeps the
exact censored likelihood, but fixes the residual distribution $F_0$ to a
single named, unimodal shape (Gaussian in Neural-CET/DKAFT, Weibull in
Ranganath et al.), which reintroduces exactly the hazard-shape ceiling
described in Section~\ref{sec:aft}; or (b) frees the shape of the
predicted distribution implicitly by abandoning the shape assumption
altogether, but at the cost of abandoning the exact likelihood too
(deepAFT, DeepR-AFT, KAN-AFT, GRAFT, DART), which reintroduces the
estimation fragility described above.} No prior deep AFT model combines
a genuinely free, learned residual distribution with the exact censored
likelihood of Eq.~\eqref{eq:censored-loglik}.

\section{Generative and flow-based survival models: the direct lineage}
\label{sec:generative}

This section covers the family of methods closest, in spirit, to this
research's own approach: models that try to learn a genuinely flexible,
shape-free residual distribution \emph{while still fitting it by the
exact censored likelihood}. To understand the two closest prior works in
detail, we first need to explain, from first principles, what a
\textbf{normalizing flow} is -- a concept from the general deep
generative modelling literature that an ML/DL reader may already know in
the abstract, but which is worth re-deriving carefully here because its
specific mathematical properties (closed-form invertibility, in
particular) are the crux of this research's technical contribution.

\subsection{Normalizing flows, from scratch}
\label{sec:flows-from-scratch}

Start from the standard \textbf{change-of-variables} formula for a
continuous random variable. Suppose $Z$ has a known, simple distribution
-- a \textbf{base distribution} -- with density $p_0(z)$ (a standard
normal, or, as in this research, a standard logistic distribution). Let
$g$ be a function that is \textbf{strictly monotone} (either always
increasing or always decreasing) and therefore \textbf{invertible}, and
define a new random variable $U = g^{-1}(Z)$, i.e., $Z = g(U)$. Then the
density of $U$ is obtained by the familiar one-dimensional
change-of-variables rule from probability theory,
\begin{equation}
  p_U(u) \;=\; p_0\bigl(g(u)\bigr) \cdot \bigl|g'(u)\bigr|,
  \label{eq:change-of-variables}
\end{equation}
where $|g'(u)|$ -- the absolute value of the derivative of $g$, the
one-dimensional Jacobian -- accounts for how much $g$ locally stretches
or compresses the probability mass as it maps $u$ to $z$. A
\textbf{normalizing flow} is exactly this construction, but with $g$
itself built to be flexible -- typically the composition of several
simple invertible transformations, often with the transformation's
detailed shape controlled by extra parameters, which in the
\emph{conditional} case (the case relevant to survival analysis) are
themselves produced by a neural network taking covariates $x$ as input.
This lets $p_U$ be, in principle, an arbitrarily complex, multimodal, or
oddly-shaped distribution, built up from a simple base distribution
$p_0$ via a learned, flexible, but always invertible transformation $g$.

\begin{intuition}
Why go to the trouble of insisting $g$ be invertible, rather than just
using any neural network to map a simple noise variable to a complicated
output (as, for instance, a GAN's generator does)? Because
Eq.~\eqref{eq:change-of-variables} gives you an \emph{exact, evaluable
density formula} for $U$ -- and an exact density formula is exactly what
Eq.~\eqref{eq:censored-loglik}, the exact censored likelihood, needs in
order to be computed at all. A GAN's generator is not (in general)
invertible, so there is no way to write down $p_U(u)$ for a specific
value $u$; you can only \emph{sample} from the model, not evaluate its
likelihood at a specific point. A normalizing flow, precisely because
$g$ is invertible and its Jacobian is tracked, gives you both: you can
sample (map a draw from $p_0$ forward through $g^{-1}$), \emph{and} you
can evaluate the exact density at any specific point (map that point
backward through $g$ and apply Eq.~\eqref{eq:change-of-variables}). This
is the entire reason normalizing flows are the natural tool for building
a flexible residual distribution for the AFT model while still
respecting the exact-likelihood discipline of
Section~\ref{sec:censored-likelihood}.
\end{intuition}

There is a second, separate question beyond whether $g$ is invertible
\emph{in principle}: whether $g^{-1}$ can be computed \emph{in closed
form} (an explicit formula) or only by running a \emph{numerical
algorithm} every time it is needed. This distinction -- closed-form
inversion versus numerical inversion -- is precisely where the two
closest prior works to this research, described next, differ from this
research's own model.

\subsection{Mixture-based approaches}
Before turning to true normalizing flows, note that the simplest way to
build a flexible, shape-free density while keeping the exact likelihood
is a \textbf{mixture model}: a weighted sum of $K$ simple component
densities (e.g., $K$ Weibull distributions), with a neural network
producing the per-subject mixture weights and component parameters from
$x$. Deep Survival Machines \citep{nagpal2022dsm} -- "DSM" -- is the
standard, well-established example of this approach, and it is one of
the main baseline models this research compares against throughout.
Structurally, however, a mixture of $K$ unimodal components is still
built from a \emph{finite sum} of simple pieces: with few components it
genuinely cannot represent a bathtub or multimodal hazard; with many
components, the number of mixture weights that must be estimated under
censoring grows, and $K$ itself becomes a fragile hyperparameter to
choose -- flexible only in a discrete, architecturally pre-committed
sense, rather than a smoothly, continuously learned one.

\subsection{Ausset et al.\ (2021): continuous normalizing flows via a neural ODE}
\label{sec:ausset}

\citet{ausset2021flows} (elaborated in the accompanying doctoral thesis,
\citealp{ausset2021thesis}) is the closest prior published work to this
research, and understanding exactly how it works -- and where it falls
short of what this research delivers -- requires one more piece of
machinery: a \textbf{continuous normalizing flow}, built from a
\textbf{neural ordinary differential equation}.

Instead of building $g$ as a composition of a fixed number of discrete
transformation "layers" (the usual normalizing-flow construction),
imagine letting the transformation happen \emph{continuously}, governed
by a differential equation:
\begin{equation}
  \frac{dz(\tau)}{d\tau} \;=\; f_\theta\bigl(z(\tau), \tau, x\bigr), \qquad \tau \in [0, 1],
  \label{eq:node}
\end{equation}
where $f_\theta$ is a neural network (the "velocity field") and $\tau$ is
an artificial "flow time" running from $0$ to $1$ (not to be confused
with the survival time $t$). Starting from $z(0) = u$ (a residual value)
and integrating this differential equation forward to $\tau = 1$ gives
$z(1) = g(u)$: the transformation $g$ is defined \emph{implicitly}, as
"wherever this differential equation ends up after one unit of flow
time," rather than by an explicit formula. This is exactly the idea
Ausset et al.\ apply to survival analysis, training the resulting flow by
the exact right-censored likelihood, and remarking, in passing, that the
construction "can be seen as a generalization of the AFT model."

\begin{intuition}
What does it cost to work with a flow defined this way? Because $g$ has
no explicit formula, \emph{every single time} the model needs a density
value, a survival probability, a sample, or a quantile, it must
\textbf{numerically solve the differential equation} in
Eq.~\eqref{eq:node} from scratch -- using a numerical ODE solver that
takes many small steps and accumulates a small amount of discretization
error at each one. This is not a one-time setup cost paid at training
time; it recurs at \emph{every} evaluation, including every single
forward and backward pass during training itself. There is no way to
write down $S(t\mid x)$ or a quantile as an explicit formula you could,
for instance, differentiate by hand or evaluate on a pocket calculator --
only "the number you get by running the solver."
\end{intuition}

Beyond this structural cost, a close reading of the Ausset thesis reveals
three further, specific gaps relative to this research's design. First,
the thesis's own review of alternative flow constructions explicitly
considers, and sets aside in favour of the continuous formulation, the
\emph{discrete}, layer-based flows this research uses instead -- so the
route this research takes was a known, deliberately declined alternative
in that prior work, not an oversight. Second, the empirical evaluation in
that work is comparatively narrow: a single multimodal synthetic
scenario, at one fixed 80\% censoring level, with Harrell's C
(Section~\ref{sec:evaluation}) as the \emph{only} reported metric -- no
hazard-shape recovery error measured against a known ground truth, no
calibration analysis of any kind, and no controlled study of how
performance changes as censoring proportion or censoring mechanism
varies. Third, and directly relevant to interpretability: although the
authors describe their model as "a generalization of the AFT," the
fitted model has no explicit $\mu(x)/\sigma(x)$ decomposition of the kind
in Eq.~\eqref{eq:aft}, and no time-ratio quantity
(Eq.~\eqref{eq:time-ratio}) is ever extracted or validated anywhere in
the thesis -- the AFT connection remains a one-sentence framing remark,
not a feature of the fitted model that a practitioner could actually use.

\subsection{Li \& Cai (2026): discrete monotone flows with no closed-form inverse}
\label{sec:licai}

The second closest work is concurrent with, and at the time of writing
still unpublished relative to, this research: \citet{licai2026}, under
review for STAI-X 2026. Where Ausset et al.\ build their flow
\emph{continuously} via a differential equation, Li \& Cai build a
\textbf{discrete monotone flow} directly: a function $g$ constructed as a
sum of scaled and shifted $\tanh$ functions (a "tanh-sum" flow), each
individual term monotone increasing, so the sum is guaranteed monotone
increasing too, and therefore invertible in principle. Because the
transform is built from an explicit formula rather than an implicit ODE
solution, the \emph{forward} direction -- computing $g(u)$, and hence the
density and survival function via
Eq.~\eqref{eq:change-of-variables} -- is available in closed form,
and the model is trained by the exact right-censored likelihood, augmented
with a distribution-matching (Soft--Nelson--Aalen Wasserstein)
regularization term and an extension addressing informative censoring
via a shared latent variable.

The catch is the \emph{reverse} direction. A sum of $\tanh$ functions
does not, in general, have a closed-form algebraic inverse -- there is no
formula you can write down for $g^{-1}(z)$ the way there is, for
instance, for a linear function or the Weibull's own quantile function.
To invert it, Li \& Cai's approach (as reviewed in the manuscript under
this project's pre-registered novelty sweep) uses \textbf{bisection
search}: a simple, general-purpose root-finding algorithm that
repeatedly halves a bracketing interval known to contain the answer,
narrowing in on the true value of $g^{-1}(z)$ one binary digit of
precision at a time, over roughly 40 iterations to reach acceptable
accuracy.

\begin{intuition}
Why does this matter? \emph{Sampling} from the model -- generating a
synthetic event time, which many downstream uses require, such as
uncertainty quantification via many simulated draws, or generating a
predicted median or other quantile of survival time for an individual
patient -- requires applying $g^{-1}$, since sampling always goes from
the simple base distribution to the data space. Under Li \& Cai's
construction, \emph{every single sample or quantile query requires
running a fresh 40-step bisection search}, an iterative numerical
procedure, rather than evaluating one closed-form expression. This is a
smaller, more contained cost than solving a full differential equation
(Section~\ref{sec:ausset}), but it is the same fundamental issue in a
different guise: \textbf{exactness in one direction (forward, toward the
likelihood) does not imply exactness in the other (backward, toward
samples and quantiles)}, and Li \& Cai's construction, like Ausset et
al.'s, is exact in only one of the two directions. Separately, because
the transform's parameters (the tanh terms' locations and scales) are
general functions of the covariates $x$ rather than entering through the
specific location-scale form of Eq.~\eqref{eq:aft}, the model has no
acceleration-factor interpretation of the kind Section~\ref{sec:aft}
described -- the map is flexible, but it is not of AFT type.
\end{intuition}

The manuscript's reported simulation study, as reviewed, uses only
standard parametric distribution families for the underlying
data-generating process (no bathtub, multimodal, or crossing-hazard
scenario is included), reports no calibration metric of the kind
described in Section~\ref{sec:evaluation}, and does not include a
controlled sweep over censoring proportion. The distribution-matching
regularizer and the informative-censoring extension are genuine,
valuable technical contributions in their own right, and this research
cites them as complementary rather than competing ideas -- but the
specific combination this research is built around (bidirectional
closed-form exactness \emph{and} an explicit AFT decomposition \emph{and}
a systematic hazard-shape/calibration study across censoring regimes)
remains open in this concurrent manuscript, as reviewed.

\subsection{A parallel statistical vocabulary: transformation models}
A separate statistical literature -- \textbf{deep conditional
transformation models} \citep{baumann2021transformation} and
\textbf{DRIFT} \citep{kolb2024drift} -- arrives at essentially the same
underlying object (a flexible, invertible, monotone transformation of a
latent variable, fit by a censored likelihood) starting from
distributional regression theory rather than from the normalizing-flow
literature. These are, mathematically, close relatives of the flow-based
models above viewed through different vocabulary and motivation, and a
further recent entrant, \citet{yuan2025transformation}, develops a
related interval-censored, partially-linear transformation model fit by
a monotone-spline sieve maximum likelihood combined with a deep neural
covariate map; none of these papers is cast in the AFT location-scale
form of Eq.~\eqref{eq:aft} or evaluated for hazard-shape recovery across
a controlled censoring sweep.

\subsection{Other, more distant approaches}
For completeness: \citet{miscouridou2018} was an early attempt (2018) at
applying an invertible transform of a Weibull base distribution to
survival data, but inside a latent-variable model fit by
\textbf{variational inference} (optimizing an evidence lower bound, an
approximation to the true likelihood rather than the exact likelihood
itself). \citet{chapfuwa2018adversarial} and the related survival-function
matching work \citep{chapfuwa2021sfm} generate event times using an
\textbf{adversarial (GAN-style)} mechanism, which imposes essentially no
shape restriction at all on the generated distribution but, as the
intuition box in Section~\ref{sec:flows-from-scratch} explained, provides
no tractable likelihood and no closed-form survival function, and is
prone to the training instability GANs are generally known for. Finally,
\citet{gharari2023copula} tackle the harder problem of
\emph{dependent} censoring using deep copula models -- explicitly out of
the scope of this research, as noted in Section~\ref{sec:censoring-types},
but the relevant reference point for anyone wishing to extend this line
of work in that direction.

\section{This research's method: FlowSurv-AFT, in full mathematical detail}
\label{sec:model}

With the whole landscape in view, this section explains exactly what
this research does, and how it resolves the tension identified at the
end of Section~\ref{sec:aft}: keep the AFT's closed-form tractability
and its time-ratio interpretability, but remove the restriction that the
residual distribution $F_0$ be a fixed, named, shape-restricted family --
\emph{and} do so with a transform that is exactly invertible \emph{in
both directions}, unlike either of the two closest prior works just
described.

\subsection{The model, stated}
The model keeps the AFT's location-scale decomposition of
Eq.~\eqref{eq:aft} exactly, but with both the location \emph{and} the
scale given by a shared neural network (an "encoder"), and with the
residual $\eps$'s distribution given by a flexible, learned, conditional
transform of a fixed, simple base distribution:
\begin{equation}
  \log T \;=\; \mu(x) + \sigma(x)\,\eps, \qquad
  \eps \;=\; g^{-1}\bigl(Z\,;\, h(x)\bigr), \qquad
  Z \sim \mathrm{Logistic}(0, 1),
  \label{eq:flowsurv}
\end{equation}
where $\mu(x)$ and $\sigma(x) > 0$ are the AFT's usual location and scale
functions (here, outputs of a neural network rather than a linear
function of $x$), $h(x)$ is a further, separate output of the same
network supplying the parameters that control the shape of the
transform $g$, and $g(\cdot\,; h(x))$ is a \textbf{conditional
rational-quadratic spline (RQS)} transform -- the specific choice of
flexible, invertible function this research uses, explained next.

\subsection{Rational-quadratic spline transforms: the technical heart of the model}
\label{sec:rqs}

The central design decision -- the one that gives this research's model
its defining advantage over both of the closest prior works
(Sections~\ref{sec:ausset}--\ref{sec:licai}) -- is the specific choice of
$g$. Rather than an implicit ODE solution (Ausset et al.) or a sum of
$\tanh$ functions with no closed-form inverse (Li \& Cai), this research
uses a \textbf{rational-quadratic spline}, following
\citet{durkan2019nsf}'s construction from the general normalizing-flow
literature, applied here to a survival residual for the first time (to
our knowledge; see Section~\ref{sec:generative} for the full novelty
argument).

Here is how it is built, piece by piece. First, choose a bounded interval
$[-B, B]$ on the residual axis (outside this interval, $g$ is fixed to be
the identity map, $g(u) = u$, so the transform only does interesting work
in a controlled, bounded region -- the "tails" behave simply and
predictably). Divide $[-B, B]$ into $K$ bins using $K+1$ \textbf{knot}
points, both on the input ($u$) axis and, correspondingly, on the output
($z$) axis, chosen (via the conditioning network's output $h(x)$) so that
the map from input knots to output knots is monotone increasing. Within
each bin, rather than connecting the two knot points with a straight
line (which would only allow $g$ to be piecewise-\emph{linear}, quite
restrictive) or an ordinary cubic polynomial (which does not, in
general, guarantee monotonicity), the transform uses a
\textbf{rational-quadratic function} -- the ratio of two quadratic
polynomials in $u$ -- with its two free coefficients chosen (again from
$h(x)$) so that the function's derivative at both bin endpoints matches a
specified, positive target value. A short calculus argument
\citep{durkan2019nsf} shows that as long as those target derivatives are
positive, the resulting rational-quadratic piece is \emph{guaranteed}
monotone increasing across the whole bin -- which is exactly the property
needed for $g$ to be invertible at all.

\begin{intuition}
Why go to the trouble of a rational-quadratic function specifically,
rather than something simpler? Because within each bin, both the forward
map $z = g(u)$ \emph{and} its inverse $u = g^{-1}(z)$ can be written down
as explicit, closed-form algebraic formulas -- solving a
rational-quadratic equation for its root is a standard, closed-form
computation (essentially an application of the quadratic formula), not
an iterative numerical procedure. And because the transform is applied
\emph{bin by bin} to bounded pieces, the derivative $g'(u)$ needed for
the Jacobian term in Eq.~\eqref{eq:change-of-variables} is likewise an
explicit formula, not something requiring automatic differentiation
through an unrolled numerical solver. \textbf{This is the single
technical fact that makes this research's model bidirectionally
exact where the two closest prior works are each exact in only one
direction}: Ausset et al.'s continuous flow needs an ODE solver for
\emph{both} directions; Li \& Cai's discrete flow is closed-form forward
but needs 40-step bisection backward; this research's spline flow is
closed-form in \emph{both} directions, with no numerical solver anywhere
in the model.
\end{intuition}

Composing several such spline "blocks" (the model uses between 1 and 3,
chosen by validation) gives the full transform $g(\cdot\,; h(x))$, and
the change-of-variables formula from
Eq.~\eqref{eq:change-of-variables} -- now conditional on $x$ through both
$h(x)$'s control of $g$ and the outer AFT location-scale transformation
of Eq.~\eqref{eq:flowsurv} -- delivers all four of the model's core
outputs in closed form:
\begin{align}
  S(t \mid x) &= 1 - F_0\bigl(g(u(t); h(x))\bigr), \\
  f(t \mid x) &= f_0\bigl(g(u(t); h(x))\bigr)\cdot\bigl|g'(u(t); h(x))\bigr| \big/ \bigl(t\,\sigma(x)\bigr), \\
  h(t \mid x) &= f(t\mid x) \big/ S(t \mid x), \\
  Q(q \mid x) &= \exp\Bigl\{\mu(x) + \sigma(x)\, g^{-1}\bigl(F_0^{-1}(q);\, h(x)\bigr)\Bigr\},
\end{align}
where $u(t) = (\log t - \mu(x))/\sigma(x)$, and $F_0$, $f_0$ are the CDF
and density of the standard logistic base distribution -- both of which
are themselves elementary, closed-form functions (the logistic sigmoid
and its derivative). \textbf{Sampling} is likewise one pass: draw $Z$
from the base logistic distribution, apply $g^{-1}(\cdot\,; h(x))$, and
exponentiate through the AFT location-scale transform -- no iteration,
no solver, no rejection sampling.

\subsection{A principled starting point: the log-logistic warm start}
One further design choice matters for both optimization stability and
for how the model's flexibility is judged experimentally. The spline
transform's parameters are initialized so that, at the very start of
training, $g$ is \emph{exactly} the identity map ($g(u) = u$ for every
$u$, achieved by initializing the conditioning network's final layer to
output zero). Substituting the identity for $g$ in
Eq.~\eqref{eq:flowsurv} collapses the model exactly to
$\log T = \mu(x) + \sigma(x) Z$ with $Z$ logistic -- in other words,
\textbf{at initialization, this research's deep flow model is
identical to a log-logistic AFT model} (one of the classical models in
Section~\ref{sec:aft}), and only departs from it as training finds
support in the data for additional residual flexibility. This matters
for the experimental design described in Section~\ref{sec:study}: it
guarantees that on data genuinely generated by a classical AFT-shaped
mechanism, the flexible model has a principled, built-in route to
recovering essentially the classical answer, rather than being forced to
"discover" a simple shape from scratch inside a much larger, generic
function class.

\subsection{Fitting the model}
The complete set of parameters -- everything controlling $\mu(x)$,
$\sigma(x)$, and $h(x)$, i.e., the entire neural encoder together with
the spline transform's conditioning -- is fit in a single pass by
directly maximizing the exact right-censored log-likelihood of
Eq.~\eqref{eq:censored-loglik}, with $f$ and $S$ substituted from the
closed-form expressions above. No partial likelihood
(Section~\ref{sec:cox}), no discretization
(Section~\ref{sec:deephit}), no inverse-probability weighting or
iterative imputation (Section~\ref{sec:deepaft}), and no approximate
lower bound in place of the true likelihood
(Section~\ref{sec:generative}) is used anywhere in fitting the model.

\subsection{What this buys, concretely}
Putting the pieces of this section together, FlowSurv-AFT delivers, in
one model: an event-time distribution whose \emph{shape} is learned from
data rather than restricted to a named family (so it can, in principle,
represent bathtub and multimodal hazards, unlike any classical AFT model
or the Neural-CET/DKAFT ablation family); every one of density, survival,
hazard, quantile, and sample available as a closed-form computation with
no numerical solver in either direction (unlike both Ausset et al.\ and
Li \& Cai); an explicit $\mu(x)/\sigma(x)$ decomposition supporting the
same time-ratio interpretation as the classical AFT
(Eq.~\eqref{eq:time-ratio}), which none of the flexible generative models
in Section~\ref{sec:generative} provide; and estimation by the same
honest, exact censored likelihood the classical AFT itself uses, with no
substitute criterion standing in for it.

\section{Why this is different: the full comparison, in one place}
\label{sec:comparison}

Table~\ref{tab:comparison-plain} restates, in the plain-language terms
this document has built up, the comparison that
Section~\ref{sec:generative} developed model by model. Read each row as
answering: "does this family of models have this property, \emph{as a
structural fact about its functional form and estimation method} --
regardless of how much data or compute is thrown at it?"

\begin{table}[h!]
\centering
\footnotesize
\setlength{\tabcolsep}{4pt}
\renewcommand{\arraystretch}{1.3}
\begin{tabular}{p{3.3cm}p{2.5cm}p{3.1cm}p{2.3cm}p{2.7cm}}
\toprule
\textbf{Model family} & \textbf{Can represent bathtub/multimodal hazards} & \textbf{Every output (density, survival, quantile, sample) closed-form, no solver} & \textbf{Has a time-ratio (AFT) interpretation} & \textbf{Fit by the exact censored likelihood, no substitute} \\
\midrule
Cox / deep Cox (DeepSurv, Cox-Time) & Yes (no shape restriction) & No (step-function baseline) & No & No (partial likelihood) \\
Weibull / log-normal AFT & No (fixed named shape) & Yes & Yes & Yes \\
Royston--Parmar spline & Yes & Yes & No & Yes \\
DeepHit / discrete-time models & Yes, within grid resolution & No (grid steps only) & No & Yes, but only over a discretized grid \\
deepAFT / DeepR-AFT / KAN-AFT / GRAFT / DART & Yes (shape unrestricted) & No (no density at all) & Yes & No (rank loss / imputation / IPCW) \\
Neural-CET / DKAFT & No (fixed Gaussian shape) & Yes & Yes & Yes \\
Deep Survival Machines (mixture) & Only approximately, via $K$ components & Yes & No & Yes \\
Ausset et al.\ (2021), continuous flow & Yes & No (ODE solver needed, both directions) & No & Yes \\
Li \& Cai (2026), discrete flow & Yes & Forward only (backward needs 40-step search) & No & Yes (forward) \\
\textbf{This research: FlowSurv-AFT} & \textbf{Yes} & \textbf{Yes, both directions, no solver} & \textbf{Yes} & \textbf{Yes} \\
\bottomrule
\end{tabular}
\caption{The full comparison in plain terms. This research's model is the
only row with a "Yes" in every column.}
\label{tab:comparison-plain}
\end{table}

\section{What this research actually does: the study design}
\label{sec:study}

Building the model is only half the contribution; the other half is
\emph{demonstrating}, with evidence rather than assertion, when and why
it should be preferred. This section explains the experimental design in
plain terms.

\subsection{A simulation study built around known ground truth}
The central methodological advantage of a \emph{simulation} study is that
the true hazard shape, the true survival function, and the true density
are all \emph{known exactly}, because the researcher wrote down the
data-generating mechanism -- unlike real clinical or reliability data,
where the truth is never directly observable and a model's error cannot
be measured against a known target, only inferred indirectly. This
research generates data under several distinct hazard-shape scenarios --
including a standard Weibull scenario (where classical AFT models should
do well, by construction), and bathtub, multimodal, and covariate-crossing
scenarios (the shapes Section~\ref{sec:math} argued a classical AFT model
\emph{cannot} represent, no matter the sample size) -- and, for each
scenario, systematically varies sample size, the proportion of censored
observations (from 0\% up to a heavy 80\%), and the censoring mechanism
(Type~I versus Type~III, Section~\ref{sec:censoring-types}), fitting
every method under comparison to every combination.

Because the true hazard is known in simulation, the study can compute a
\textbf{hazard-recovery error}: a direct numerical distance between each
model's estimated hazard curve and the true one, at each covariate
profile, on a fine time grid -- the single most direct test of whether a
model actually recovers the shape it was designed to recover, as opposed
to indirect proxies like concordance. This, together with
distributional-fidelity measures (comparing estimated and true cumulative
distribution functions directly) and the calibration metrics from
Section~\ref{sec:evaluation}, is what this research's central hypotheses
(restated in plain language below) are tested against.

\subsection{The four hypotheses, restated in plain language}
\begin{itemize}
  \item \textbf{H1a (flexibility -- shape and location perturbations,
    confirmatory).} On hazard shapes that are still, underneath,
    variations in the \emph{shape or location} of a single residual
    distribution -- a multimodal (two-hump) hazard, or a hazard whose
    location $\mu(x)$ moves nonlinearly with covariates -- FlowSurv-AFT
    should recover the true hazard more accurately than every classical,
    machine-learning, and deep baseline tested, including the
    mixture-based comparator (DSM).
  \item \textbf{H1b (flexibility -- mechanism-misaligned shapes,
    descriptive only).} A pilot run partway through this research
    surfaced an important nuance, worth stating plainly rather than
    papering over: some bathtub- or crossing-hazard scenarios are
    generated by a mechanism that does not sit naturally inside the AFT
    picture at all. Recall from Eq.~\eqref{eq:aft} that the AFT
    decomposition lets covariates \emph{shift and stretch the time axis}
    ($\mu(x)$, $\sigma(x)$) and, in this research's version, also
    \emph{reshape the residual distribution} through the flow's
    conditioning $h(x)$ -- but if the true mechanism instead multiplies
    the \emph{hazard itself} by a covariate factor at every instant (the
    Cox-style form $h(t\mid x) = h_0(t)\cdot\exp(\beta^\top x)$, which is
    how one of the bathtub scenarios in this study's simulation design
    is actually generated), that covariate effect has to be represented
    entirely through the harder, less naturally-optimized shape pathway
    $h(x)$ rather than the location pathway the model is warm-started
    toward (Section~\ref{sec:model}). A reduced-budget tuning pass, two
    increases in the model's spline capacity, and an experiment that
    deliberately broke the identity warm start all failed, independently
    and consistently, to close a gap against DSM observed on exactly
    these scenarios -- evidence that this is a genuine property of the
    mechanism, not simply a case of the model being undertrained or
    badly configured. For these specific scenarios, then, this research
    does not claim FlowSurv-AFT beats DSM; the comparison is reported
    honestly, without a predicted direction, rather than forced to fit
    the original blanket hypothesis. This is, if anything, a more useful
    scientific answer than a clean sweep would have been: it says
    precisely \emph{which} kinds of hazard-generating mechanisms an
    AFT-shaped flow model helps with, and which ones still call for a
    different modelling approach altogether.
  \item \textbf{H2 (calibration under censoring).} Models trained by the
    exact censored likelihood should hold up better, in calibration
    terms, as censoring gets heavier, than models trained by discretized
    or partial-likelihood surrogates -- and the size of this advantage
    should \emph{grow} as censoring proportion increases, for the
    estimator-fragility reasons discussed in Section~\ref{sec:deepaft}.
  \item \textbf{H3 (no penalty for flexibility).} On data that genuinely
    is Weibull-shaped -- exactly the setting a classical AFT model is
    designed for -- FlowSurv-AFT should not do meaningfully worse than
    the classical Weibull AFT itself, thanks to the log-logistic warm
    start described in Section~\ref{sec:model}: the extra flexibility
    should cost nothing when the data does not need it.
  \item \textbf{H4 (interpretability preserved).} Where a classical AFT
    model is already an adequate description of real data, this
    research's learned time ratios (Eq.~\eqref{eq:time-ratio}, computed
    from the fitted $\mu(x)$) should agree closely with the classical
    model's own time ratios -- evidence that the flexible model's
    acceleration-factor estimates can be trusted -- while, on covariates
    where the classical linear model is \emph{not} adequate, the fitted
    model should reveal genuinely nonlinear acceleration effects a
    linear AFT cannot represent at all.
\end{itemize}

\subsection{Real-data validation}
Beyond the simulation study, the model is also fit to five public
clinical benchmark datasets covering a range of real censoring
proportions and covariate structures, to check that the simulation
findings transfer to real, uncontrolled data. Consistent with the
evaluation-methodology discussion in Section~\ref{sec:evaluation}, the
real-data claims this research makes are deliberately restricted to
calibration and distributional quality rather than concordance,
because -- per \citet{burk2024neutral} -- concordance improvements on
standard tabular clinical data are not, in general, expected to be
available to any method, flexible or otherwise; the honest place to look
for an advantage is whether the predicted \emph{distributions} are
better calibrated, especially on the most heavily censored of the five
datasets.

\section{Summary}
\label{sec:summary}

Survival analysis exists because some outcomes are only partially
observed, and every model in this field is, at bottom, a choice of
functional form for the event-time distribution together with a choice
of how to fit that form using both fully- and partially-observed data
correctly (Sections~\ref{sec:why}--\ref{sec:math}). Reviewing the field's
history model by model (Sections~\ref{sec:classical}--\ref{sec:generative})
reveals a consistent pattern: every existing approach gains flexibility,
tractability, interpretability, or exact-likelihood estimation by giving
up at least one of the other three. This research's model,
\textbf{FlowSurv-AFT}, is constructed specifically to occupy the gap left
by that pattern: it keeps the classical accelerated failure time model's
closed-form tractability and its interpretable time-ratio effect measure,
replaces the classical model's restrictive fixed residual distribution
with a rational-quadratic spline transform that is exactly invertible in
\emph{both} directions -- unlike either of the two closest prior flow-based
models, each of which is exact in only one direction -- and fits the
whole model by the same honest, exact censored likelihood the classical
model itself uses, with no approximation or substitute criterion
anywhere. The accompanying simulation and real-data study is designed to
show, with direct evidence rather than assertion, exactly when this
combination pays off, and to be explicit -- following the best current
evaluation-methodology practice in the field -- about the one thing it is
not designed to claim: that it ranks patients better than a fifty-year-old
regression model, because the best available evidence says that claim is
not there to be made on data of this kind in the first place.

\bibliographystyle{apalike}
\bibliography{references}

\end{document}
